\begin{frame}{How it is argued, not just whether it is true}
  \adminikicker{innovation \textperiodcentered\ rhetorical pass}
  \vspace{2mm}

  \adminilead{A claim can be accurate and still be deployed to mislead. That
    reading runs beside fact-checking, never inside it.}

  \vspace{3mm}
  \begin{itemize}
    \item It names the technique---emotional pressure, loaded language, faulty
          reasoning, selective evidence, conspiracy framing.
    \item Every signal must quote a verbatim transcript line or it is dropped;
          the prompt names the technique, never the position.
  \end{itemize}
  \note{%
    This runs beside fact-checking rather than inside it: a true statistic with
    its base rate removed, a sponsorship read as editorial, a contested
    political question framed as settled. Each signal must quote a verbatim
    transcript line or it is dropped, and the prompt names the technique, never
    the position. It applies the same threshold whatever the politics: if it
    would not flag the mirror-image passage from the opposing side, it does not
    flag this one. An empty list is the normal, correct answer for an ordinary
    video. With citation verification, this is the other half of the 30\% for
    idea and innovation.%
  }
\end{frame}
